\title{X-LoRA: Mixture of Low-Rank Adapter Experts, a Flexible Framework for Large Language Models with Applications in Protein Mechanics and Molecular Design}

\begin{abstract}
We introduce a mixture of expert methodology for fine-tuning large language models utilizing a deep, layer-wise, token-level technique grounded in low-rank adaptation (LoRA). Beginning with a collection of pre-trained LoRA adapters, our gating mechanism leverages hidden states to dynamically combine adapted layers, enabling the X-LoRA model to tap into diverse capabilities and form novel deep layer-wise combinations to address various tasks. This design draws inspiration from biological concepts of universality and diversity, where neural network components are repurposed across different hierarchical levels. Consequently, the X-LoRA framework can be seamlessly applied to any existing large language model (LLM) without requiring modifications to its fundamental architecture. We develop a specialized X-LoRA variant that provides scientific functionalities, including forward and inverse analysis tasks and improved reasoning abilities, with an emphasis on biomaterial analysis, protein mechanics, and design. The significance of this work lies in providing models that are easily extensible and adaptable, rich in domain expertise, and capable of integrating multiple knowledge domains. Incorporating experts in biology, mathematics, reasoning, bio-inspired materials, mechanics and materials science, chemistry, protein biophysics, mechanics, and quantum mechanics-based molecular properties, we perform a suite of physics-centered case studies. These include assessments of knowledge retrieval, protein mechanics forward and inverse challenges, protein design, adversarial agentic modeling with ontological knowledge graph construction, and molecular design. The model not only produces quantitative predictions of protein nanomechanical properties and quantum mechanical molecular characteristics but also interprets outcomes and accurately identifies plausible mechanisms explaining distinct molecular behaviors.
\end{abstract}

\section{Introduction}
Large language models (LLMs)~\cite{Vaswani2017AttentionNeedc,Touvron2023LlamaModels,OpenAI2023GPT-4Report,Chowdhery2022PaLM:Pathways,Jiang2023Mistral7B,Gunasekar2023TextbooksNeed} have attracted considerable attention, including efforts to create specialized models that excel in particular tasks, reasoning processes, or scientific areas~\cite{Bubeck2023SparksGPT-4,Buehler2023MechGPTModalities,Nejjar2023LLMsAnalysis,Buehler2023GenerativeDesign,Luu2023BioinspiredLLM:Materials,Luu2023GenerativeSolvents,Buehler2023MeLMProblemsb,Ge2023OpenAGI:Experts}. Nonetheless, training such specialized models can be expensive, particularly when a broad range of capabilities is required. Approaches like low-rank adapters (LoRA) \cite{hu2021lora} have been suggested as a more resource-efficient alternative, although these adaptations typically target narrower knowledge domains. The fundamental idea behind LoRA is to employ low-rank matrices added to the original full-scale matrix, with only these low-rank components being trainable within the model. Because training is confined to the adapter layers, these models generally retain the knowledge acquired during the base model’s pre-training, while enhancing applicability to particular tasks and reducing computational demands. Due to the efficiency of training LoRA adapters, it becomes feasible to create a wide array of specialized models using this method.

Although LoRA adapters offer an efficient route to develop specialized model functionalities, a key open question remains about how to combine multiple such adapters into unified models that provide a cohesive and potentially enhanced set of abilities. Indeed, the combination of LLMs has emerged as an intriguing research direction~\cite{Kim2023SOLARUp-Scaling}. Additional experimental approaches, such as spherically interpolated model parameters (SLERP) and related techniques~\cite{Arcee-ai/mergekit:Models.}, have demonstrated promising outcomes and suggest the potential for combining models to produce novel capabilities. These mixed models are typically constructed {\it a priori} based on specific algorithms, and while they have shown encouraging results on certain benchmarks, further investigation is required to better understand the rational design of mixed models and the mechanisms through which particular features arise.

We propose that dynamically mixing parameters represents a more general technique, as it enables the combined model to investigate novel parameter combinations during inference.

Similar methodologies have employed mixture of experts (MoE) frameworks~\cite{Jacobs1991AdaptiveExperts,Hampshire1992TheRecognition,Jordan1994HierarchicalAlgorithm}, including the recently introduced Mixtral LLM~\cite{Jiang2024MixtralExperts}. Nevertheless, these methods often involve high computational costs and substantial resource requirements even at inference time. In this work, we introduce a computationally efficient mixture-of-expert strategy that utilizes a collection of separate LoRA adapters, each of which can be trained independently with ease. Our method features a highly granular and adaptable mechanism, drawing inspiration from a biological model that reutilizes neural network components to build a hierarchy of interacting models, ranging from the base model to agent-based systems (Fig.~\ref{fig:hierarchical_concept}). To demonstrate its effectiveness quantitatively, we concentrate our investigation on applying this approach to science-oriented LLMs, specifically targeting protein mechanics challenges within the broader context of bio-inspired materials analysis and design.

We present an algorithm capable of dynamically mixing adapters at the token-level state, with fine granularity allowing mixing at the level of individual layers. This technique, termed X-LoRA, provides access to diverse capabilities essential for scientific applications of multimodal language modeling~\cite{Hu2023DeepScience,Buehler2023AMetamaterials,Buehler2022ProteinNetworks,}.

\subsection{Fundamental concepts of X-LoRA}

We begin with a concise review of the principles underlying the development of LoRA adapters. The foundational idea of low-rank adaptation (LoRA)~\cite{hu2021lora} posits that weight updates possess a low “intrinsic dimension,” which is exploited by fixing the original weight matrix $W_0 \in \mathbb{R}^{d \times k}$ and limiting modifications to a low-rank decomposition

\begin{equation}
W_0 + \Delta W_0 = W_0 + BA 
\end{equation}

where $B \in \mathbb{R}^{d \times r}$ and $A \in \mathbb{R}^{r \times d}$ with rank $r \ll {\rm min}(d,k)$. During training with LoRA~\cite{hu2021lora}, the pretrained weights $W_0$ remain fixed, and only the matrices $A$ and $B$ contain trainable parameters. The forward pass of a LoRA layer can be represented by the following equation:

\begin{equation}
    h = W_0x + \Delta W_0x = W_0x + BAx
\end{equation}

In practice, LoRA adapters include a fixed scalar weight $\alpha$, which is applied to the decomposition matrices, resulting in

\begin{equation}
    h = W_0x + BAx \cdot \alpha
\end{equation}

Building upon this, X-LoRA is founded on the concept of employing a collection of LoRA adapters, each possessing distinct capabilities, as demonstrated in previous science-oriented LLM research (e.g., \cite{Buehler2023GenerativeDesign}). Rather than statically mixing or integrating models, we introduce a dynamic gating mechanism that adjusts the scaling of each LoRA adapter with fine granularity at the token and layer levels to enable mixing deep within the model. Specifically, for a LoRA adapter $i$, the original parameter $\alpha_i$ is multiplied by a scaling factor $\lambda_i$ to produce the new scaling value $\alpha^*_i$ for that adapter:

\begin{equation}
    \alpha^*_i = \alpha_i \cdot \lambda_i
\end{equation}

The scaling factor $\lambda_i$ is generated by an X-LoRA scaling head, which leverages the model's hidden states and constitutes the trainable element in the X-LoRA architecture. Because the scaling head can utilize complex representations, it can be efficiently implemented as a simple feed-forward neural network with one or a few layers (in the X-LoRA implementation, users can define the configuration of this neural network).

The scaling computation is performed individually for each token in the sequence, thus providing a fine-grained level of control as each LoRA adapter operating at specific layers of the LLM is scaled accordingly.

From this point forward, we designate this strategy of gating the adapter components in this manner as scaling. Additional details about the method are provided in the Materials and Methods section.

\subsection{Paper outline}

The structure of this paper is as follows. Initially, we describe the methodology and training strategy, including the pertinent datasets, which lead to the creation of an X-LoRA model endowed with specialized capabilities in the physical sciences, with a particular emphasis on biomaterials such as proteins. Following this, we demonstrate a series of experiments where we employ the X-LoRA for various tasks, including question answering, conversational and agentic modeling, protein design and analysis, among others. We perform an in-depth analysis of scaling trends and validate the method by comparing it with molecular modeling and other physical data and techniques.

\section{Results and discussion}

The development process of the X-LoRA model involves the following stages:

\begin{enumerate}
    \item Training the foundational base LLM (completed in prior work) \item Individually training a collection of adapters to cultivate expertise in several domains (these represent distinct or overlapping skill and knowledge areas) \item Training the integrated X-LoRA model
\end{enumerate}

Our experimental work begins with training a series of LoRA adapters. We create a set of nine adapters, each fine-tuned to acquire specific expertise, based on the Zephyr-7B-$\beta$ model~\cite{HuggingFaceH4/zephyr-7b-betaFace}, which itself is built upon the Mistral-7B model~\cite{Jiang2023Mistral7B}:

\begin{enumerate}
    \item Bioinspired materials \item Chain-of-thought (CoT) and reasoning \item Chemistry \item Mathematics \item Physics \item Biology \item Mechanics and materials \item Platypus: Logic and reasoning \item Protein mechanics tasks (generative sequence-to-property and inverse capabilities)
\end{enumerate}

These particular focus areas and domains of expertise were selected because our ultimate goal is to develop a STEM-oriented, materials science-centered LLM. Nonetheless, we hypothesize that our methodology can be broadly extended to various fields, especially to enable the creation of modalities that facilitate links between disparate knowledge domains~\cite{Buehler2023MechGPTModalities,Buehler2023MeLMProblemsb}. 

Following the creation of the LoRA adapters, we proceed to train the X-LoRA model. This training utilizes a subset of the combined collection of training samples originally employed for the individual LoRA adapters (refer to Materials and Methods for detailed descriptions of the dataset, implementation specifics, and related information). 

The X-LoRA framework serves as an experimental tool to evaluate its capabilities across multiple scenarios. We carry out a range of experiments, including complex question answering, deploying the generative model for protein design, performing analysis and reasoning, as well as developing advanced interaction strategies through adversarial agentic modeling. Throughout these investigations, we compare the base model’s performance against that of the X-LoRA model, monitoring how the scaling head allocates weights across different layers of the various experts during answer generation or development. 

\subsection{Question answering and observed X-LoRA layer-wise scaling weights}

Our initial analysis focuses on contrasting the X-LoRA model with the original foundational model. Both models are posed the following question:

\begin{quote}
\small\texttt{What material behavior (stiffening vs softening) causes supersonic crack speeds in mode II cracks? Think step by step.}
\end{quote}

The base model’s reply is:

\begin{quote}
\small\texttt{Supersonic crack speeds, also known as Rayleigh waves, occur in mode II cracks due to the stiffening material behavior of the material near the crack tip. In mode II fracture, the loading is shear-dominated and the crack propagates perpendicular to the direction of applied stress. Here's a step-by-step explanation:

As the crack tip nears the material, the stresses around it escalate rapidly. This creates a zone of intense stress concentration known as the stress singularity.

Within this stress singularity, the material undergoes elastic deformation that causes an increase in stiffness, referred to as stiffening behavior. This stiffening arises because the material adjacent to the crack tip becomes increasingly constrained as the crack advances [...]}
\end{quote}

Drawing on our own expertise~\cite{Buehler2008AtomisticFailure}, this response is not only inaccurate but also unnecessarily verbose. Conversely, the X-LoRA model answers succinctly:

\begin{quote}
\small\texttt{The material behavior that causes supersonic crack speeds in mode II cracks is stiffening.}
\end{quote}

This response is accurate, clear, and concise~\cite{Buehler2008AtomisticFailure}. 

The distinction between the base model and the X-LoRA model becomes even more pronounced for specialized protein analysis and design tasks. The given task is:

\begin{quote}
\small\texttt{CalculateForceHistory< M S K S N T Y R M L V L L E D D T K I N K E D E K F L K G K P G K M H E F V D E L I L P F N V D E L D E L N T W F D K F D A E I C I P Y V E I L K E S G M K > }
\end{quote}

The base foundation model responds with:

\begin{quote}
\small\texttt{I'm sorry but the input provided is not a valid mathematical expression. It appears to be a sequence of characters that do not form a calculable function. Please provide a valid mathematical expression for me to calculate.

However, I can tell you that the sequence of characters you provided is actually a series of function and variable names that may be part of a larger program or script in a specific programming language. Without context, it's impossible to determine the intended calculation.

If you provide more information, such as the programming language and the context in which this expression appears, I may be able to provide more insight into its meaning and possible calculations. }
\end{quote}

The X-LoRA model provides the following response, accurately performing the design task of producing the force-deformation curve linked to this protein:

\begin{quote}
\small\texttt{[0.013,0.066,0.078,0.045,0.072,0.126,0.127,0.107,0.119, 
0.074,0.106,0.098,0.103,0.128,
0.077, 0.069,0.131,0.108,0.144,0.174, 0.166,0.183,0.185,0.217,0.175,0.230,0.187,0.224, 0.267,0.301,0.315]}
\end{quote}

Figure~\ref{fig:QA_weights_sample_0} presents results from this question answering along with the observed X-LoRA scaling weights. The two queries—one targeting mechanics and materials (Fig.~\ref{fig:QA_weights_sample_0}a) and the other focusing on protein mechanics (Fig.~\ref{fig:QA_weights_sample_0}b)—exhibit distinct patterns. Notably, we observe that a complex arrangement of scaling values is employed in each instance, implying that the X-LoRA model utilizes a heterogeneous combination of adapters distributed across layers. Additionally, the heatmap in Fig.~\ref{fig:QA_weights_sample_0}a displays a wider range of activations spanning multiple experts and layers, including some regions of intense activation. In contrast, the heatmap in Fig.~\ref{fig:QA_weights_sample_0}b reveals a more focused activation pattern, with fewer widespread high values. We infer that this might suggest the gating function on the right operates more selectively in activating experts.

In Fig.~\ref{fig:QA_weights_sample_0}, the bright line-like regions correspond to paths of strong activation, indicating that the scaling function selectively enhances certain experts for particular layers. The observed patterns imply that the importance is not evenly spread across all experts and layers; rather, it is quite sparse and selective. Such sparsity aligns with findings from other mixtures-of-experts models, where only a limited subset of experts is activated for any given input to specialize in distinct aspects of the problem space.

By comparing the activation maps across various tasks, both the sparsity and the patterns may provide valuable insights into which experts the model depends on for different input types or during different stages of neural network processing (we will discuss how these maps evolve dynamically when addressing complex tasks later). 

To further investigate the model’s behavior, we now examine a wider range of questions to understand how the X-LoRA model engages different experts given distinct prompts. Our particular focus is on whether and how the X-LoRA model combines multiple adapters for questions that span overlapping skills or domains. Fig.~\ref{fig:QA_weights} presents results from question answering alongside the observed X-LoRA scaling weights, with the questions detailed in Table~\ref{tab:table_qa}. In Table~\ref{tab:table_qa}, we provide both the X-LoRA-generated responses and those from the base model, Zephyr-7B-$\beta$. The differences between these two are pronounced.

To highlight this clearly, incorrect or questionable responses are marked with a \redhighlight{red highlight}. The comparison demonstrates that X-LoRA delivers significantly improved accuracy and more concise answers. 

As anticipated, we observe a complex interplay of adapters, often involving the activation of multiple dominant LoRA experts. For example, Fig.~\ref{fig:QA_weights}d shows that a question regarding the mechanics of pine cones strongly activates the mechanics/materials adapter along with the CoT/reasoning and bioinspired adapters. Although X-LoRA selects the correct answer letter, the explanation is not completely accurate since pine cones open in dry, not wet, conditions. However, a subsequent follow-up question probing the release mechanism yields the correct answer. 

In a protein-related query, Fig.~\ref{fig:QA_weights}f demonstrates adapter mixing in response to a question about the mechanics of Bombyx mori silk, engaging a combination of bioinspired, mechanics/materials, and biology adapters. For highly specialized tasks such as protein design (Fig.~\ref{fig:QA_weights}a), we observe primarily the activation of a single expert.

A more detailed examination of the heatmap of scalings shown in Fig.~\ref{fig:QA_weights}a reveals that besides the clearly prominent importance of the protein mechanics expert, there exist multiple bands where specific experts receive significant emphasis over several layers. This is especially evident for the mechanics/materials and reasoning experts. These bands are discontinuous and exhibit a fluctuating pattern, which may indicate particular conditions or features that these experts manage effectively. Notably, the reasoning expert (positioned immediately to the left of the protein mechanics expert) is distinguished by having the highest activation in several layers, particularly in the upper layers (approximately layers 100 to 140). We hypothesize that this may suggest a crucial role in the model's final decision-making or output, although further investigation is required to confirm this. Additionally, there are distinct but localized strong activations for the physics and biology experts at certain layers, possibly reflecting a specialization of these experts for features or representations within those segments.

Moreover, Fig.~\ref{fig:QA_token_history} presents a token-level history of the aggregated scalings across all layers. This analysis highlights histories for two examples: Fig.~\ref{fig:QA_token_history}a relates to a task concerning pine cone seed release, while Fig.~\ref{fig:QA_token_history}b depicts the outcome of a sequence of protein mechanics tasks. The protein mechanics task included a series of prompt-answer interactions: initially two property calculations, followed by a generative task aimed at designing a new protein towards the conclusion. Although expert usage is generally consistent, noticeable variations occur in the deployment of the LoRA experts throughout the process. Specifically, the protein mechanics expert is heavily engaged in the final phases of generation when the new protein design is underway.

Table~\ref{tab:table_qa} presents various use cases spanning different domains along with a comparison to the performance of the base model operating independently. We explore several of these cases in greater detail.

For example, we present a question that combines protein science and biology with logical reasoning, focusing on protein expression patterns. X-LoRA provides the correct answer. The base model's response, however, contains some errors and logical flaws when evaluating the expression of marker proteins in each cell type. The base model accurately recognizes that fibroblasts express Protein B and do not express Protein A due to the mutual exclusivity between Protein A and Protein B. Nevertheless, it fails to explicitly mention that fibroblasts also express Protein C. Since there is no restriction against co-expressing Protein C with Protein B, and fibroblasts have no limitations on Protein C expression, it logically follows that they also express Protein C. Regarding neurons, the base model correctly starts by indicating that neurons express Protein A and not Protein B, in line with the mutual exclusivity rule. However, its discussion on whether neurons might express Protein C or both Protein C and Protein B is unclear and contains an error. Given that neurons express Protein A and not Protein B, the only valid option per the rules is that neurons express Protein C. There is no circumstance in which neurons could express Protein B, contrary to the initial analysis. Hence, neurons express Protein A and Protein C, but not Protein B. Concerning epithelial cells, the base model correctly states that epithelial cells do not express Protein A but do express both Protein B and Protein C. However, its explanation is complicated and includes incorrect claims about mutual exclusivity between Protein A and Protein C, which was not part of the original constraints. The straightforward reasoning is that epithelial cells do not express Protein A and must therefore express the other two proteins, Protein B and Protein C, since the initial instructions only prohibit the co-expression of Protein A and Protein B.

X-LoRA, on the other hand, accurately determines the marker proteins expressed by each cell type by utilizing the given observations and logical inferences. For fibroblasts, X-LoRA correctly concludes that these cells cannot express Protein A because the presence of Protein A excludes the expression of Protein B. Since fibroblasts do express Protein B, and considering that each cell type expresses a unique combination of three marker proteins (an incorrect assumption introduced, as the original statement did not specify the number of proteins expressed but only that the combination is unique), the inference that fibroblasts also express Protein C is valid based on the provided information (excluding the part about the number of proteins, which appears to be a misinterpretation). Regarding neurons, the model correctly identifies them as expressing Protein A and Protein C. The rationale is valid, recognizing that Protein B cannot be co-expressed with Protein A, and since neurons express Protein A and another marker protein (not Protein B), they must express Protein C. Lastly, for epithelial cells, the X-LoRA model accurately concludes that they express Protein B and Protein C because epithelial cells do not express Protein A, and by elimination, given that they express two marker proteins, these must be Protein B and Protein C.

In the question concerning pine cone release in humid versus dry environments, the base model provides a convoluted and factually inaccurate response. X-LoRA, however, delivers a clear and precise answer with well-articulated reasoning. Additional related instances include the question about the exceptional mechanical properties of Bombyx mori silk, where the base model incorrectly claims that silk fibers conduct electricity. The base model also fails to solve several other domain-specific problems, such as the question regarding the role of hyperelasticity in maximum fracture speeds, for which X-LoRA offers a succinct and accurate response while the base model does not provide an answer. In the question about the strong adhesion of gecko feet, X-LoRA correctly highlights the significance of setae and nanoscale effects, whereas the base model erroneously attributes the process to special-purpose glue proteins. These and other examples presented in the table and throughout the paper strongly support the superior capabilities of X-LoRA compared to the base model.

Figure~\ref{fig:perf_comp_bioinspiredexam} presents a comparison of knowledge recall evaluation experiments conducted with the bio-inspired knowledge test set introduced in \cite{Luu2023BioinspiredLLM:Materials}. It is evident that X-LoRA achieves the highest performance despite being a considerably smaller model than the BioinspiredLLM or Llama-BioLLM models (7B parameters for X-LoRA compared to 13B parameters in the others). Additional comparisons include other recently released LLMs such as Orca-13B and Llama-13b-chat.

The figure also contrasts X-LoRA with its base model, Zephyr-7B-$\beta$, demonstrating that X-LoRA exhibits significantly enhanced performance. These comparisons suggest that robust performance can be attained with X-LoRA even when employing smaller base models.

Figure~\ref{fig:image_synth} illustrates an application in image synthesis, where X-LoRA is used to generate a prompt for other models, such as DALL-E 3, as shown here.

\begin{quote}
\small\texttt{Create a prompt that I can use for image generation that accurately describes the microstructure of a hierarchical composite.

Explicitly describe the geometry.}
\end{quote}

The comparison (Figure~\ref{fig:image_synth}(b)-(c)) clearly demonstrates that the more concise prompt leads to a significantly more accurate image generated by DALL-E 3~\cite{DALLECard}. Notably, the prompt created by Zephyr-7B-$\beta$ does not precisely adhere to the given instruction and introduces additional design elements such as fibers and nanoparticles, which were not mentioned in the original prompt focused exclusively on the microstructure of a hierarchical composite. The prompt generated by X-LoRA is shorter and more targeted, resulting in a far superior outcome.

\subsection{Generative solutions to protein analysis}

We now turn to specific protein generation and analysis tasks. These abilities emerge in the X-LoRA model through the protein mechanics adapter and are further enhanced by integration with the model’s existing knowledge of biology, bio-inspired materials, reasoning, and logical deduction. This section focuses on analyzing protein-related tasks independently, evaluating the quantitative performance of this functionality.

Fig.~\ref{fig:protein_force_history_prediction} presents outcomes from the protein task aimed at predicting force-deformation behavior based on sequences. Fig.~\ref{fig:protein_force_energy_prediction} illustrates the overall accuracy in forecasting unfolding force and unfolding energy from the sequence, comparing actual values with model predictions (yielding an $R_2=0.85$). In general, the model demonstrates excellent performance, and as will be further examined in the following section, it also performs exceptionally well on the inverse problem of design and cycle-consistent evaluation\cite{Lu2023GenerativePropertiesb}. Nevertheless, unlike previous studies that employed highly specialized GPT-type models for these tasks, our X-LoRA model integrates these specialized functionalities with a variety of other deep specializations introduced through multiple fine-tuned adaptations.

\subsection{Protein design and analysis}

We now broaden the scope of test cases to include more comprehensive applications that engage multiple capability domains. Specifically, the experiments discussed here demonstrate how the X-LoRA model leverages a diverse set of skills and how agentic modeling can generate even more complex and thorough responses. We start by addressing a protein design task, where the model is tasked with creating a novel protein tailored to achieve a specified target force-deformation profile. These force-deformation responses have been investigated in previous research~\cite{Ni2023ForceGen:Model} and represent measurements of force as a function of deformation during protein stretching at both ends, characteristic of a classical protein unfolding experiment~\cite{Lu1998UnfoldingSimulation.}.

Fig.~\ref{fig:design_protein} illustrates the application of the generative protein task. We engineer a protein to exhibit a specific force-deformation response (training data and boundary conditions derived from the original molecular dynamics (MD) simulations of protein pulling, see~\cite{Ni2023ForceGen:Model}), then evaluate the predicted sequence to determine if it achieves the intended target properties (Fig.~\ref{fig:design_protein}a). This cycle-consistent evaluation is crucial for verifying that the inverse design aligns with the forward prediction.

In Fig.~\ref{fig:design_protein}b, the resulting predicted sequence is shown: \texttt{MSKSNTYRMLVLLEDDTKINKEDEKFLKGKPGKMHEFVDELILPFNVDELDELNTWFDKFDAEICIPYVEILKESGMK}. Additionally, Fig.~\ref{fig:design_protein}c presents a folded structure of this protein, generated using AlphaFold 2~\cite{Jumper2021HighlyAlphaFoldb}. To understand the relationship of the designed protein to known sequences, Fig.~\ref{fig:design_protein}d displays its comparison via a Basic Local Alignment Search Tool (BLAST) Tree~\cite{Altschul1990BasicTool}. Notably, the designed protein shows connections to hypothetical proteins investigated in prior studies but which remain structurally unresolved.

Leveraging agentic modeling and integrating multiple X-LoRA models within a multi-agent framework, we further examine the predicted sequence through our X-LoRA model. This approach capitalizes on the model’s advanced capability to not only interpret protein mechanics but also incorporate interdisciplinary knowledge. Using the sequence predicted above as input, we request a more thorough analysis and suggestions for sequence variations. We begin with the following query:

\begin{quote}
\small\texttt{I have identified this amino acid sequence: MSKSNTYRMLVLLEDDTKINKEDEKFLKGKPGKMHEFVDEL ILPFNVDELDELNTWFDKFDAEICIPYVEILKESGMK. 

How can I use it? 

Provide logical steps, discuss specifically the engineering of the problem. }
\end{quote}

Through multiple agentic interactions employing adversarial strategies (details in Materials and Methods), we encourage the model to explore more nuanced and sophisticated aspects of the topic. The summary of this interaction is presented below as identified by the X-LoRA model (the full transcript is available in the Supplementary Material, SI Conv Example 1):

\begin{quote}
\small\texttt{The discussion focuses on the potential uses of a particular amino acid sequence in the field of protein engineering. This sequence features a combination of hydrophobic and hydrophilic residues, implying that it could establish a stable conformation and interact with other molecules. To design a protein incorporating this sequence, computational approaches can be employed to forecast its secondary and tertiary structures; however, challenges include the intricate nature of protein folding and the high computational expense. Enhancements can be achieved by creating more efficient algorithms and integrating new experimental findings. It is essential to validate computational models to guarantee their precision and practical applicability. Methods such as cross-validation, regularization, early stopping, ensemble techniques, and transfer learning are useful to avoid overfitting and underfitting. Furthermore, interpreting and explaining these models can support practical implementation.}
\end{quote}

We also request the model to generate a summary of key points, resulting in:

\begin{quote}
\small\texttt{\begin{itemize}
    \item The amino acid sequence MSKSNTYRMLVLLEDDTKINKEDEKFLKGKPGKMHEFVDELI LPFNVDELDELNTWFDKFDAEICIPYVEILKESGMK holds promise for protein engineering due to its blend of hydrophobic and hydrophilic residues. Computational tools can predict the protein’s secondary and tertiary structures, though obstacles include the complexity of protein folding and computational demands.
    \item Validating computational models is vital to confirm their accuracy and practical value. Techniques like cross-validation, regularization, early stopping, ensemble methods, and transfer learning can mitigate risks of overfitting and underfitting.
    \item Model interpretation and explanation further aid in guiding real-world applications.
\end{itemize}
}
\end{quote}

Ultimately, the primary conclusion is:

\begin{quote}
\small\texttt{The key insight is that the amino acid sequence MSKSNTYRMLVLLEDDTKINKEDEKFLKGKPGKMHEFVDELILPFNVDELDELNTWFDKFDAEICIPYVEILKESGMK holds promise for protein engineering applications due to its combination of hydrophobic and hydrophilic residues. To design a protein with this sequence, computational approaches can be employed to predict its secondary and tertiary structures, though challenges include the complexity inherent in protein folding and the computational resources required. Ensuring the accuracy and practical relevance of these computational models requires rigorous validation, with methods such as cross-validation, regularization, early stopping, ensemble approaches, and transfer learning helping to mitigate overfitting and underfitting. Additionally, interpreting and explaining the models aids in guiding practical implementations.}
\end{quote}

We can also delve deeper into the particular steps for sequence engineering. To facilitate this, we prompt the model to offer more targeted design recommendations aimed at enhancing the protein’s stability:

\begin{quote}
\small\texttt{How might I modify the following protein sequence to improve its stability: MSKSNTYRMLVLLEDDTKINKEDEKFLKGKPGKMHEFVDELILPFNVDELDELNTWFDKFDAEICIPYVEILKESGMK? }
\end{quote}

The interaction summary is as follows (complete dialogue available in Supplementary Material, SI Conv Example 2):

\begin{quote}
\small\texttt{The discussion centers on strategies to increase protein stability by incorporating hydrophobic amino acids and forming disulfide bonds. The biologist notes that hydrophobic residues generally enhance stability, and that disulfide bonds can further strengthen the protein structure. The engineer requests concrete examples of hydrophobic amino acids and guidance on the number of disulfide bonds to introduce. The biologist cites examples of proteins successfully stabilized through these modifications, such as insulin, bovine pancreatic trypsin inhibitor, and green fluorescent protein. The conversation also addresses potential drawbacks of adding hydrophobic amino acids or disulfide bonds, including altered protein function, heightened structural complexity, and risks of aggregation. The biologist advises that these challenges can be managed via careful amino acid selection, computational modeling, and experimental validation.}
\end{quote}

The list of main points:

\begin{quote}
\small\texttt{\begin{itemize}
\item To enhance the stability of a protein sequence, one can incorporate hydrophobic amino acids and disulfide bonds. 

Hydrophobic amino acids generally contribute to protein stability, while disulfide bonds also serve to increase stability.

Proteins like insulin, bovine pancreatic trypsin inhibitor, and green fluorescent protein have been effectively stabilized through these modifications.
\item Possible drawbacks or challenges of adding hydrophobic amino acids or disulfide bonds include alterations in protein function, greater complexity, and a risk of aggregation.
\item These challenges can be mitigated by carefully selecting amino acids, utilizing computational tools, and conducting experimental verification.
\end{itemize}
}
\end{quote}

Ultimately, the main conclusion is:

\begin{quote}
\small\texttt{The most crucial insight from the discussion is that to stabilize the protein sequence MSKSNTYRMLVLLEDDTKINKEDEKFLKGKPGKMHEFVDELILPFNVDELDELNTWFDKFDAEICIPYVEILKESGMK, the introduction of hydrophobic amino acids and disulfide bonds is effective. This directly addresses the initial query.}
\end{quote}

From the viewpoint of protein engineering, these recommendations are logical since they can indeed result in more stable protein designs. Concrete examples and methodologies are provided for researchers to investigate further.

As another illustration, we analyze the following three sequences. We first instruct the model to evaluate the stability of each, then ask it to determine the strongest candidate and explain why that particular protein is probably the most stable. Additionally, we request the model to outline the steps to synthesize it in the laboratory. The three sequences (manually generated by modifying one from the test set) are:

\begin{quote}
\small\texttt{
\begin{itemize}
    \item LNTWFDKFDAEICIPYVEILKESGMK \item AITWFDKFDAEICIPYVEALKIAAMK \item AAAAAAAAKFDAAEICIIAAAAAAAAKIAAMK
\end{itemize}
}
\end{quote}

This scenario evaluates if the model can integrate multiple skills and adapt dynamically throughout the conversation—that is, to fluidly switch between different abilities such as understanding protein mechanics, performing rational analysis, generating novel hypotheses, and so forth.

The conversation is summarized in Figure~\ref{fig:conversation_evolve}, alongside a representation of how the scalings evolve throughout the dialogue (the analysis includes both the deep layer-wise adapter scaling and an aggregated analysis to provide a measure of the effective utilization of the various experts). The findings clearly indicate that the protein task is initially highly significant, progressively transitioning to a stronger emphasis on the bioinspired adapter and the biology adapter. By the conclusion of the conversation, the biology adapter emerges as the most dominant.

To evaluate the predictions and assess the quality of the responses, the three proteins involved in this task are illustrated in Figure~\ref{fig:protein_visuals} (proteins folded using Alpha Fold 2~\cite{Jumper2021HighlyAlphaFoldb}, via ColabFold \cite{Mirdita2022ColabFold:All}). We observe that the most stable structure in the center exhibits the most organized secondary structure, consistent with the idea that it produces the highest unfolding force, as forecasted by the X-LoRA model in Figure~\ref{fig:conversation_evolve}. The other two structures are less ordered; the first is primarily helical with some turns, and the last is partially unstructured and features a kink in its geometry. These characteristics account for the lowest unfolding force predicted by the model.

By examining the visual depictions of the folded conformations, we confirm that the central structure, which is the most stable, exhibits the most well-organized secondary structure. This observation aligns with the concept that it produces the greatest unfolding force, as predicted by the model. The other two structures are found to be less organized: the first is predominantly helical with some turns, while the last one is partially disordered and contains a bend in its geometry. We propose that these characteristics probably account for the lowest unfolding force observed, consistent with the predictions of the X-LoRA model.

In a second instance, we focus on the energetic aspects of proteins, examining the unfolding energy as the protein is stretched from its native state to a fully extended conformation~\cite{Ni2023ForceGen:Model}. We again analyze three sequences (the first two are the same as those considered above, while the third sequence is longer, representing a combination of the second sequence at the ends and the first sequence in the middle):

\begin{quote}
\small\texttt{\begin{itemize}
    \item LNTWFDKFDAEICIPYVEILKESGMK \item AITWFDKFDAEICIPYVEALKIAAMK \item AITWFDKFDAEICIPYVEALKIAAMKLNTWFDKFDAEICIPYVEILKESGMKAITWFDKFDAEICIPYVEALKIAAMK
\end{itemize}
}
\end{quote}

Initially, we request the model to calculate the unfolding energy for each sequence, followed by an analysis of the results: Based on the amino acid sequences, we ask for an explanation as to why one protein is likely the most stable. Finally, we inquire about the potential functions of the protein identified as the most stable.

Figure~\ref{fig:MJB_20} presents a transcript of a dialogue that leverages combinations of various experts and synthesizes knowledge across these domains, applied here to the unfolding energy of three proteins. By the conclusion of the conversation, the CoT/reasoning adapter emerges as the dominant component, as the X-LoRA model answers a query requiring reasoning over the results and forecasting probable protein structural characteristics along with potential protein functions. Specifically, the mechanistic inquiry regarding the basis of protein stability includes the formation of hydrophobic interactions among nonpolar amino acids as well as hydrogen bonding among polar amino acids. The model forecasts that several hydrophobic amino acids (I, F, W, Y) and hydrogen bonding amino acids (D, E, K) contribute to the protein’s stability.

Additionally, the model predicts that cysteine (C) residues present in the sequence would facilitate disulfide bond formation, further reinforcing the protein structure. It is noteworthy that the critical predictions about structural features explaining the high stability are accurately identified, as confirmed by structural analyses shown in Fig.~\ref{fig:MJB_21}. The analysis in this figure corroborates the existence of all these features. Moreover, it becomes evident why the other two proteins, which are considerably shorter segments with variable helical domain stability, exhibit much lower unfolding energies.

\subsection{Adversarial agentic modeling to connect distinct scholarly disciplines and knowledge yielding ontological knowledge graph generation}

Given the extensive cross-domain knowledge embedded in our X-LoRA model, we can utilize it to explore links between diverse ideas, knowledge repositories, and fields of expertise. In one experiment, we pose two queries to the model, each framed similarly but emphasizing different perspectives. Subsequently, we generate triplets to create an ontological knowledge graph that refines the responses into more structured outputs~\cite{Giesa2012CategoryDesign,Spivak2011ReoccurringAnalogies}. The resultant graph delivers an integrated comprehension of the produced insights and visually represents the connections between concepts in a clear and mechanistic fashion.

We employ two opposing agents to address the question. Agent 1 is labeled as a "critical physicist," while Agent 2 assumes the role of a "philosopher." Agent 1 initiates the inquiry and is responsible for probing and formulating further questions in response to each answer provided by Agent 2. As the dialogue progresses, the discussion delves deeper into the topic, revealing various facets. The question is deliberately selected to span multiple disciplinary domains to assess the model's capability to synthesize ideas and concepts.

The question posed is:

\begin{quote}
\small\texttt{Consider an abstract representation of a classical musical composition as a graph of interacting notes, chords, melodies, rhythms and so on.

What happens if we would apply pressure, up to a point of failure, on this abstract concept of music?

Provide logical steps, mathematical reasoning, and discuss specifically the physics of the problem.}
\end{quote}

The entire conversation is provided in the Supplementary Material (see, SI Conv Example 3). For conciseness, we present only a summary and the key takeaway here, capturing the core aspects:

\begin{quote}
\small\texttt{The discussion explores the effect of applying pressure to an abstract representation of a classical musical composition, treating it as a physical system composed of interacting notes, chords, melodies, and rhythms. Pressure can be represented quantitatively via stress, with the failure point defined when the stress attains a critical threshold. To validate the mathematical model’s predictions, one can employ a physical model of the abstract musical composition and compare experimental observations against theoretical expectations. The model's limitations include simplified interaction representations, static assumptions, uniform distribution, and a linear correlation between stress and pressure. Enhancements to the model’s precision and validity would involve integrating more detailed and realistic interactions, accounting for dynamic phenomena, non-uniform distributions, and incorporating non-linear terms.

[...]

The principal insight is that applying pressure to the abstract musical composition can be interpreted as modifying the force fields among the interacting notes, chords, melodies, and rhythms. Increasing pressure raises the system’s internal stress, and at a critical threshold, the system will fail and collapse under its own constraints. This explanation addresses the original question by offering a physical rationale for how pressure influences the abstract notion of music and how this effect can be quantified through stress.}
\end{quote}

Although the generated text offers valuable insights, creating an ontological representation of the data can enhance interpretability~\cite{Buehler2023MechGPTModalities}. For graph generation, we merge the full content of both conversations to build a knowledge graph. An example graph is illustrated in Figure~\ref{fig:graph}. The constructed graph is partitioned into communities using the Girvan-Newman algorithm~\cite{Girvan2002CommunityNetworks}, resulting in a total of 12 communities. The figure displays both the complete graph (Figure~\ref{fig:graph}a) and detailed views with node labels (Figures~\ref{fig:graph}b-c).

\subsection{Creation of X-LoRA-Gemma Integrating Protein, Chemical, Bio-inspired, and Materials Mechanics Features}
\label{gemma-section}

To demonstrate that our proposed methodology is compatible with other foundational models featuring different architectures, we trained a second X-LoRA model, this time built upon the Gemma-7B-it model~\cite{Google/gemma-7b-itFace}. In this instance, we employ a set of four LoRA adapters, which are defined as follows:

\begin{enumerate}
    \item Bioinspired materials \item Mechanics and materials \item Protein mechanics tasks (including generative sequence-to-property and inverse design capabilities) \item Quantum mechanics-based molecular properties QM9 (incorporating generative SMILES-to-property and inverse design functionalities; see Table~\ref{tab:table_QM9} for definitions and summaries of all properties computed within that dataset)~\cite{Ruddigkeit2012EnumerationGDB-17,Ramakrishnan2015ElectronicSpace}
\end{enumerate}

This X-LoRA-Gemma model is developed in a manner similar to the original X-LoRA model but is designed to incorporate distinct functional capabilities, particularly adding the ability to predict twelve quantum mechanical properties and to design molecules that satisfy these twelve quantum mechanical properties. This new adapter is trained simultaneously with the others listed above, using the same datasets employed in the prior experiments. It is important to note that while the main body of this paper concentrates on the Zephyr-based X-LoRA model, within this section we differentiate the models by referring to the original as X-LoRA-Zephyr and the Gemma-based model as X-LoRA-Gemma.

Figure~\ref{fig:protein_force_energy_prediction} illustrates the overall performance of the X-LoRA model, highlighting predictions of unfolding force and unfolding energy derived from the sequence ($R_2 = 0.87$), as well as predictive accuracy for quantum mechanical molecular properties from the QM9 dataset ($R_2 = 0.96$). We observe excellent outcomes in both tasks, with performance on the protein design task being marginally superior to that of the original X-LoRA model.

We further showcase the capabilities of X-LoRA-Gemma through a straightforward multi-agent design task. We examine a small organic molecule, \texttt{CCCC12CC1C(=O)N2} (extracted from the QM9 dataset test set). This molecule has a bicyclic structure with a ketone (=O) and an amine (N) integrated into the ring; refer to Figure~\ref{fig:QM9_design}(a) for a depiction of this molecule. Our first query to the X-LoRA-Gemma model is:

\begin{quote}
\small\texttt{Here is a molecule in SMILES code: CCCC12CC1C(=O)N2

What changes would I need to make to make this molecule more useful for polymer design. }
\end{quote}

The reply is:

\begin{quote}
\small\texttt{To enhance the utility of this molecule for polymer design, you might consider:
\begin{enumerate}
\item  Adding a reactive functional group, like an amine or hydroxyl group, to enable further modification and crosslinking.
\item  Incorporating a longer, more flexible carbon chain to improve the molecule's capacity to form polymers with diverse structures and characteristics.
\item  Attaching a side chain that contains a functional group, such as a carboxyl group, to facilitate covalent bonding with other molecules and polymers.
\item  Embedding a cyclic structure, for example, a lactam ring, to yield a polymer that is more rigid and stable.
\item  Including various side chains to generate polymers with a broad range of properties and functionalities.
\item  Varying the number of carbons in the backbone to allow synthesis of polymers with differing molecular weights and attributes.
\item  Introducing a crosslinking site in the molecule to produce a polymer that is more stable and less reactive.
\item  Incorporating a repeating unit framework to achieve a polymer with defined molecular weight and properties.
\item  Combining several of the above strategies to develop a versatile and functional polymer for multiple applications.
\end{enumerate}}
\end{quote}

We now concentrate on the first recommendation.

\begin{quote}
\small\texttt{How can I enhance the overall reactivity of a molecule? I intend to use this molecule for constructing a polymer material.

The molecule under consideration, given in SMILES notation, is: CCCC12CC1C(=O)N2

I can adjust the following properties:

Dipole moment, mu Polarizability, alpha Highest Occupied Molecular Orbital energy, homo Lowest Unoccupied Molecular Orbital energy, lumo HOMO-LUMO gap: The energy difference between HOMO and LUMO, gap Electronic spatial extent Zero-point vibrational energy Heat capacity at constant volume Internal energy at 0 K Internal energy at 298.15 K Enthalpy at 298.15 K Free energy at 298.15 K What are the top THREE properties I should modify, and in which direction (increase or decrease) should I alter them? }
\end{quote}

The feedback from X-LoRA-Gemma is:

\begin{quote}
\small\texttt{\begin{itemize}
    \item  Decrease in HOMO-LUMO gap \item  Increase in dipole moment \item  Increase in polarizability
\end{itemize}
By altering these properties, the overall reactivity of the molecule can be enhanced, making it more appropriate for polymer material construction. }
\end{quote}

We leverage both the generative and predictive functionalities of X-LoRA-Gemma to engineer these modifications in molecular properties.

For the optimization process, we employ a two-agent framework consisting of two instances of X-LoRA-Gemma. One agent performs the generative task \texttt{GenerateMolecularProperties<...>} to produce candidate molecules. The second agent then executes the \texttt{CalculateMolecularProperties<...>} task to compute and evaluate the discrepancy between the target and predicted properties. This cycle is repeated until the maximum number of iterations is reached or the mean-squared error falls below a predefined threshold.

The outcome of this approach is illustrated in Figure~\ref{fig:QM9_design}. Panel (a) presents the target design alongside its divergence from the original molecule's properties. Panel (b) displays the top 20 molecular designs generated by the algorithm, with panel (c) showing the distribution of the mean-squared error and panel (d) plotting the mean-squared error across the generated candidates. The most successful design is \texttt{CC(C)C1=CC(=O)NC=C1}, which aligns well with the target as depicted in panel (e). Lastly, panel (f) features an optimized 3D molecular structure (optimized using UFF~\cite{Rappe1992UFFSimulations} for the molecule's geometry).

To further evaluate the molecule’s reactivity, we calculated the Gasteiger charges~\cite{Gasteiger1980IterativeCharges} (visualized in Figure~\ref{fig:QM9_design}(f)). The molecule contains an aromatic ring with a ketone (C=O) group, which is more electrophilic relative to carbon-carbon or carbon-hydrogen bonds, rendering it prone to nucleophilic attack. The combination of the nitrogen-containing aromatic ring and the ketone yields a pyridone ring structure. This ring exhibits higher reactivity than a simple aromatic ring due to the heteroatom (nitrogen) and ketone group, which provide sites for potential reactions such as nucleophilic addition at the carbonyl carbon or electrophilic substitution on the ring. The nitrogen atom within the ring, as part of the heteroaromatic system, influences the aromaticity and chemical reactivity by modifying electron density and its distribution over the ring. This effect impacts the molecule's interactions with other chemical species, positioning the nitrogen-containing ring as a locus for chemical reactions, especially those involving electrophilic attack facilitated by the electron-donating nature of the nitrogen.

The carbon atom in the ring, adjacent to both an oxygen and a nitrogen atom, carries a partial positive charge (+0.25), indicating it is somewhat deficient in electrons. This occurs because oxygen, being more electronegative, draws electron density toward itself, reducing the electron richness of the carbon. This effect is partially offset by the neighboring nitrogen, which is less electronegative than oxygen but more so than carbon. Nonetheless, the partial positive charge implies that this carbon is a probable site for nucleophilic attack, where a nucleophile (an electron-rich entity) would be attracted to this electron-poor carbon.

The nitrogen atom, possessing a partial negative charge (-0.33), suggests it is relatively electron-rich compared to its typical state. This can be attributed to its lone pair of electrons and the influence of the aromatic ring's electron system. In aromatic heterocycles, the nitrogen's lone pair can be involved in the delocalized $\pi$-electron system; however, the exact degree depends on the ring’s configuration and the electronegativity of neighboring atoms. The partial negative charge implies that nitrogen is more nucleophilic, meaning it is more prone to donate electrons, either by forming bonds with electrophiles or undergoing protonation reactions. Additional quantum mechanical studies are required to explore this further.

The carbon bearing the partial positive charge represents a reactive center for nucleophilic attacks. Other molecules may interact with this site via mechanisms where nucleophiles target electron-deficient carbons.

The hydrogen atom attached to the nitrogen, which carries a positive partial charge of 0.17, is well-suited to form hydrogen bonds with other molecules or atoms. Hydrogen bonding is a form of dipole-dipole interaction that occurs when a hydrogen atom bonded to a highly electronegative atom (such as nitrogen, oxygen, or fluorine) interacts with another electronegative atom possessing a lone pair of electrons. This characteristic enables the molecule to function as a hydrogen bond donor, an important factor affecting its solubility in water and other polar solvents, as well as its interactions within biological environments.

The ketone group positioned next to the nitrogen-containing aromatic ring imparts distinctive reactivity that can be leveraged in polymer synthesis~\cite{McQuarrieSimonPhysicalChemistry,CareySundbergAdvancedOrganicChemistry,Varghese2022BeyondApplications,Varghese2022BeyondApplications}. This compound is classified as a lactam, due to the cyclic amide structure formed by the ketone and the nitrogen within the ring. Lactams play a significant role in polymer chemistry, especially in the production of polyamides, a category of polymers widely utilized in various material applications. The ketone group can engage in multiple chemical reactions relevant to polymer chemistry, including nucleophilic addition. Such reactions can be employed to lengthen the polymer chain or to add side chains that alter the polymer’s characteristics. The nitrogen atom within the aromatic ring can notably influence the polymer’s electronic attributes and can participate in hydrogen bonding, thereby affecting the polymer's melting temperature, solubility, and mechanical behavior. N-heteroaromatic compounds are also recognized for their capacity to coordinate with metal ions, potentially enabling uses in fields like catalysis or in materials possessing electronic or photonic functionalities. This molecule could also undergo polymerization via the amide (lactam) group; polyamides synthesized in this manner are well-known for their durability, flexibility, and resistance to abrasion and chemicals, finding applications across diverse domains such as textiles (e.g., nylon) and automotive components~\cite{Varghese2022BeyondApplications,Varghese2022BeyondApplications}.

Alkyl substituents can enhance the hydrophobic nature of a molecule, so we expect that polymers synthesized from monomers containing such alkyl groups will tend to be more hydrophobic. This change influences their solubility in various solvents and their interactions with water. Such characteristics may be beneficial for applications that require water resistance or particular solubility profiles. Additionally, the incorporation of alkyl groups can modify the polymer's mechanical properties. For example, they may function as plasticizers within the polymer network, thereby increasing the material's flexibility. Depending on the length and branching of the alkyl chains, this could result in polymers that are more flexible or exhibit improved impact resistance.

\subsubsection{Additional benchmarking} We present further benchmark results for the X-LoRA-Gemma model, as illustrated in Figure~\ref{fig:perf_MMLU}, utilizing a subset of the MMLU benchmark that concentrates on chemistry, physics, and biology (fields closely related to the domains on which our model was trained)~\cite{Hendrycks2020MeasuringUnderstanding}. Our findings indicate that both X-LoRA models discussed in this work outperform their respective base models. Specifically, the X-LoRA models surpass the base models with the Zephyr model achieving 56.6\% compared to 54.6\% for the base, and the Gemma model reaching 52.4\% versus 40.6\% for the base. The highest overall accuracy is attained by the X-LoRA-Zephyr model, though the performances of both models are fairly comparable.

\section{Conclusion}

Multimodal LLMs have numerous existing and prospective applications in scientific fields, alongside a strong demand for highly specialized, domain-specific expertise. Beyond general language processing, the ability to perform computational or generative tasks—such as those illustrated here with protein mechanics—is essential~\cite{Luu2023BioinspiredLLM:Materials,Buehler2023MechGPTModalities,Buehler2023MeLMProblemsb}. X-LoRA, developed on top of open-source language models, allows for the dynamic combination of expertise and knowledge across different domains, as evidenced by various demonstrations, including Figure~\ref{fig:conversation_evolve}. This investigation revealed how the X-LoRA model autonomously and adaptively evolves its use of adapters, blending specific adapters in a sophisticated manner across the LoRA layers. Such intricate mixing of LoRA layers is consistently observed in all the examples analyzed here, such as Figures~\ref{fig:QA_weights_sample_0} and \ref{fig:QA_weights}. This phenomenon also presents an intriguing challenge to better understand the underlying processes from a physics standpoint, where one can consider an LLM as a large-scale graph-forming system that leverages complex graph structures. Consequently, the integration of adaptation experts at deep layer scales, as implemented here, modifies these mechanisms to develop effective solutions.

The algorithm introduced here extends the typical inference approach—usually a single forward pass—by incorporating two forward passes. This procedure trains the model to initially ‘consider’ the question and how it might adjust itself before generating a response. This approach embodies a basic form of ‘self-awareness,’ enabling the model to modify its own structure to optimally address a given task. Therefore, despite having a relatively modest parameter count of 7 billion, the model is capable of reasoning across a wide range of scientific domains (including biological materials, mathematics, physics, chemistry, logic, and mechanics), greatly improving its ability to produce innovative solutions. Moreover, this approach could inspire the design of alternative model architectures, such as methods that generalize the reconfiguration strategy to subsets of the model or even the entire non-adapted system. In such cases, the scaling head would dynamically adjust parts of a trained model or integrate multiple models. As demonstrated with X-LoRA, these strategies can serve as powerful tools for expanding or combining the functionalities of different models.

We conducted multiple performance evaluations, including a comparison with significantly larger models (Figure~\ref{fig:perf_comp_bioinspiredexam}), where X-LoRA demonstrated clear superiority. The comprehensive comparison presented in Table~\ref{tab:table_qa} indicated that X-LoRA consistently outperforms the base model across a range of tasks and domains of application. Additional comparisons involved quantitative evaluations of numerical prediction challenges, such as protein mechanics and quantum-mechanical molecular properties, where X-LoRA achieved high accuracy with $R_2$ values ranging from 0.85 to 0.96. For reference, these results surpassed previously published outcomes, such as protein unfolding force predictions using the MeLM model~\cite{Buehler2023MeLMProblemsb}, which reported an $R_2$ of 0.78. The ability to accurately address multiple prediction tasks, alongside inverse design problems and advanced reasoning skills, was highlighted in the molecular design example (Section~\ref{gemma-section}). This case also demonstrated that X-LoRA can be effectively applied to different base model architectures (Gemma versus Zephyr/Mistral). Future research could extend these efforts by developing X-LoRA models compatible with larger base models.

One limitation of the proposed method is the necessity of performing two forward passes, which may result in increased computational overhead: one pass to compute hidden states serving as inputs to the X-LoRA scaling head, and another with $\Lambda$ applied to the adapters. However, the initial pass is utilized for the X-LoRA scaling head, thereby leveraging inherent LLM capabilities instead of requiring separate learning through training a larger scaling head. For practical deployment in model-serving environments, distinct key-value caches should be maintained for both the scaling and forward passes. To enhance inference speed, we introduced \texttt{Mistral.rs}, an LLM serving platform implemented in Rust~\cite{matsakis2014rust} that supports X-LoRA and incorporates multiple performance optimizations. Despite this, a key benefit of our approach is its compatibility with any model without necessitating internal structural changes. We consider this a significant advantage, particularly due to its applicability to the extensive resources available within the Hugging Face ecosystem.

One benefit of the proposed method is that it offers an easy way to integrate into any existing LLM (for example, since the X-LoRA code has been made available, it should be straightforward to use with any model within the Hugging Face ecosystem). Another benefit is that the scaling head leverages the inherent capabilities of the LLM and learns sophisticated approaches for optimally combining the adapters to generate specific answers. When training X-LoRA, it is ideal to use complex question-answer pairs or conversations that enable the model to discover the most effective combinatorial strategies. We conducted several analyses related to protein design and mechanics. For example, Figure~\ref{fig:MJB_20} and Fig.~\ref{fig:MJB_21} showcase stability analyses of multiple protein sequences, including comparative assessments, reasoning based on predicted outcomes, and mechanistic explanations that offer a foundational chemical rationale for the observed behavior. As illustrated in Fig.~\ref{fig:MJB_21}, this was directly validated through a structural analysis of the protein's folded 3D conformation.

Regarding future directions, further studies could investigate a wider array of adapter experts, expanding on the initial expertise embedded in our LoRA set. The X-LoRA model surpasses the capabilities of the base model or a base model equipped with a single adapter, enabling it to address complex questions by drawing on specialized skills from the ensemble of adapters. There also appear to be interactions among the different adapters, as suggested by the intricate mixing patterns of scaling weights across layers. This phenomenon warrants additional investigation and comparison with approaches such as SLERP that combine models using alternative techniques. While we trained the X-LoRA scaling head using a subset of training data—consisting of a few hundred samples from each original training set used to create the adapters—a more targeted design of training datasets for this purpose could be developed. For example, we observed that posing multiple-choice questions tends to activate reasoning-focused adapters, since these were trained on such data. Additional prompting may be necessary to help the scaling head comprehend specific domains required to answer certain questions, potentially through specialized system messages or other methods. Overall, creating appropriate training sets is an area ripe for further research, including methods to effectively induce mixing of layer-wise scaling. Our experiments tracking these mechanisms over extended conversations demonstrate that the model can dynamically adjust scaling to optimally address different tasks, which is an encouraging capability.

We identified intriguing patterns in how experts are activated across layers, with noteworthy insights elaborated in the paper. This analysis has the potential for further extension and could provide valuable understanding of the benefits of mixing model components or groups of layers. Such exploration would be especially compelling when applying the X-LoRA approach to domains beyond protein mechanics and protein design, as demonstrated here. We reserve this for future research.

Another significant direction involves deeper investigation of agentic modeling, exemplified here by the use of two X-LoRA models engaging in adversarial interaction. Future studies might involve more than two models to enrich interactions and introduce additional functionalities, thereby advancing models into new generative territories. This strategy could also support the creation of methods that combine physics-based modeling or other validation procedures, leveraging code-writing/executing agents or agents employing function calling and other processing techniques~\cite{Ghafarollahi2024ProtAgents:Learning}.

\section{Materials and Methods}

The X-LoRA codebase and sample implementations can be accessed at \url{https://github.com/EricLBuehler/xlora}. Model weights and datasets related to the X-LoRA discussed herein, along with supplementary examples, are hosted at \url{https://huggingface.co/lamm-mit/x-lora}. The X-LoRA-Gemma model is available at \url{https://huggingface.co/lamm-mit/x-lora-gemma-7b}. The \texttt{Mistral.rs} inference engine can be found at \url{https://github.com/EricLBuehler/mistral.rs}.

\subsection{Mathematical formulation of X-LoRA} We define the $\mathbf{scalings}$ as a matrix encompassing token- and layer-level coefficients for LoRA adapters. The X-LoRA $\mathbf{scaling}$ $\mathbf{head}$ receives the hidden states from the base model and processes them through fully connected linear layers to predict scalings, followed by a softmax function along the final dimension to ensure that scalings across all experts sum to one. ReLU activation functions and dropout layers are interspersed between the linear layers in the X-LoRA scaling head. The hidden layer sizes are configurable and can be adjusted to effectively reduce dimensionality from the hidden dimension to the scaling dimension, or to implement a widening-and-narrowing scheme analogous to the feed-forward layer in transformer blocks.

The scalings form a matrix $\Lambda \in \mathbb{R}^{s \times l \times n}$, where $s$, $l$, and $n$ denote the sequence length, the number of layers, and the number of adapters, respectively. For each layer, the relevant scalings are extracted as a matrix $\lambda \in \mathbb{R}^{s \times n}$. These scalings $\lambda$ are then applied to each LoRA adapter by broadcast multiplication of the corresponding X-LoRA scaling $\lambda_i \in \mathbb{R}^{1 \times 1 \times s}$ with the input to the decomposition matrices $A$ and $B$.

Given $W_0 \in \mathbb{R}^{d \times k}$ and adapter-specific matrices $B_i \in \mathbb{R}^{d \times r_i}$ and $A_i \in \mathbb{R}^{r_i \times d}$, where the rank $r_i \ll \min(d,k)$, the forward pass of an X-LoRA layer is defined as follows:

\begin{equation}
    h = W_0x + \sum_{i=1}^{n} B_iA_i(x \cdot  \lambda_i) \alpha_i
\end{equation}

Visual representations of the scalings are provided in the main text. These were illustrated either via bar plots or through the use of the \texttt{contourf} function in Matplotlib~\cite{MatplotlibDocumentation}.

\subsection{X-LoRA implementation algorithm and code} The chosen implementation prioritizes user-friendliness and the flexibility to easily apply the method to a range of models, particularly those within the Hugging Face framework. The X-LoRA algorithm functions by:

\begin{enumerate}
    \item Keeping the base model and adapters frozen while fine-tuning performance.
    \item Combining LoRA adapters based on predictions from the X-LoRA scaling head.
\end{enumerate}

All programming is done using Python and PyTorch~\cite{Paszke2019PyTorchLibrary}. 

\subsubsection{Dual forward pass strategy for self-aware inference} Since X-LoRA necessitates computing hidden states to estimate the scaling factors, we adopt a dual forward pass method. This involves an initial pass where the model assesses the task and determines how to adjust itself, followed by the actual inference step. This approach introduces a form of `self-awareness' that enables X-LoRA to allocate extra computation time for reflection instead of an immediate response. 

We refer to the first pass as the scaling pass, and the second as the forward pass. During the scaling pass, the scalings $\Lambda$ are set to zero. We also tested setting the scalings $\Lambda$ to ${n}^{-1}$, observing comparable outcomes that suggest stable performance (this parameter is configurable in the X-LoRA package). To facilitate easy application of X-LoRA across different models, we implement a forward pass hook. This mechanism is critical in linking the scaling and forward passes and represents the key feature enabling the method’s applicability to all models.

\begin{enumerate}
    \item $\mathbf{Scaling}$ $\mathbf{pass}$: Executed by running the base model with adapters fixed at constant values. The resulting hidden states feed into the X-LoRA scaling head to produce $\Lambda$.
    \item $\mathbf{Forward}$ $\mathbf{pass}$: Uses $\Lambda$ to blend the adapters during the model's forward pass. The resulting logits constitute the output of the X-LoRA model.
\end{enumerate}

A word of caution: using the same key-value cache for both scaling and forward passes can disrupt correct input to the scaling head, as the cache would persist and accumulate over the two passes. Therefore, we disable the key-value cache during both training and inference, implementing proper management of this behavior in \texttt{Mistral.rs}.

\subsubsection{Freezing weights and enabling model segments for training} In the X-LoRA model, the base model and adapters remain frozen (meaning their weights are not updated during training), making the X-LoRA scaling head the sole component that is trainable. The code provided allows the option to unfreeze all adapters along with the X-LoRA scaling head if desired. Additionally, one might choose to train the entire model—comprising the X-LoRA scaling head, adapters, and base foundation model—either simultaneously or with varying learning rates. Furthermore, future developments based on the concepts introduced here could facilitate creating scaling functions between two or more foundation models, enabling a transparent implementation of SLERP-type modeling. This direction is left for future research.

\subsubsection{X-LoRA layers} An X-LoRA layer functions by scaling and then combining the outputs from multiple adapters. It wraps around the original LoRA layer and thus has access to the specific adapters associated with that layer.

To effectively integrate X-LoRA within any model, the X-LoRA algorithm replaces the original LoRA layers by constructing X-LoRA layers and injecting their forward pass logic into the LoRA layer. This process leverages features of Python’s object-oriented system to enable the swapping.

In brief, during the forward pass, each X-LoRA layer performs the following steps (refer to Figure \ref{fig:general_arch}):

\begin{enumerate}
    \item $\mathbf{Scaling}$: Multiply the input signal of each LoRA adapter by its corresponding scaling factor $\lambda_i$.
    \item $\mathbf{Forward}$: Execute the forward computation of the encapsulated LoRA layer.
\end{enumerate}

\subsection{Efficient Inference with \texttt{Mistral.rs} implemented in Rust} To enhance inference efficiency, we created \texttt{Mistral.rs} in Rust, a large language model serving platform that supports X-LoRA and includes multiple performance optimizations~\cite{matsakis2014rust}. Among the implemented techniques, without any specific ordering, are:

\begin{itemize}
    \item LoRA adapter weight stacking: When all $A_i$ and $B_i$ matrices share the same shape across each LoRA adapter, it is possible to stack the LoRA adapter weight matrices $A_i$ and $B_i$ along their first dimension.

Mathematically, given $W_0 \in \mathbb{R}^{d \times k}$, for adapter matrices $A_i \in \mathbb{R}^{r \times d}$ and $B_i \in \mathbb{R}^{d \times r}$ with $r_i \ll {\rm min} (d,k)$, we concatenate these matrices resulting in $A\in \mathbb{R}^{i \times r \times d}$ and $B\in \mathbb{R}^{i \times d \times r}$. The parameters $\alpha_i$ are also applied to the stacked matrices during initialization.

To implement scalings, we rearrange the dimensions so that $\lambda \in \mathbb{R}^{i \times 1 \times 1}$, after which they are applied via broadcast multiplication.

\item Non-granular scalings To minimize the frequency of running the scaling operation, we provide an option for non-granular scalings. This approach caches the scalings after $n$ iterations because the KV cache guarantees the sequence length remains one for the rest of the request. Utilizing this method significantly lowers latency by reducing the number of calls to the base model.

\item Fused Compute Unified Device Architecture (CUDA) kernels To enhance the efficiency of the forward pass, we develop fused CUDA kernels. These fused kernels allow complex operations, such as Rotary Embedding (\cite{su2021roformer}), which would typically execute sequentially, to be merged into a single kernel executed in parallel on a GPU. Our kernels are based on the vLLM implementation (\cite{kwon2023efficient}).

\item Quantization The inference engine supports quantization, enabling the use of quantized base models. Preliminary tests show that quantization up to 8 bits performs well even on models initially trained with \texttt{bfloat16}, for example.

\end{itemize}

\subsection{Datasets}

Each adapter is trained using specialized datasets, summarized in Table~\ref{tab:table_loradata} along with references to the data sources and original literature.

\subsection{Training strategy} Training proceeds in phases. We start with the Zephyr-7B-$\beta$ model~\cite{HuggingFaceH4/zephyr-7b-betaFace}, which builds upon the Mistral-7B model~\cite{Jiang2023Mistral7B}, and train a sequence of adapters using the datasets outlined in the previous section. The Zephyr-7B-$\beta$ tokenizer is employed.

The initial eight adapters are trained with question-answer pairs as provided by the respective datasets (see Table~\ref{tab:table_loradata}).

For protein mechanics tasks, training is conducted similarly but using particular forward and inverse instruction sets. The bidirectional instruction tasks include:

\begin{quote}
\tiny {\texttt{CalculateForce< G E E C D C G S P ....> [0.310]\\
CalculateEnergy< G E E C D C G S P ....> [0.220]\\ 
CalculateForceHistory< G E E C D C G S P ....> [0.004,0.034,0.125,0.142,0.159,0.102,0.079,0.073,0.131, ...]\\
GenerateForce<0.310>\,[ G E E C D C G S P ....]\\
GenerateEnergy<0.220>\,[ G E E C D C G S P ....]\\
GenerateForceHistory<0.004,0.034,0.125,0.142,0.159,0.102,0.079,0.073,0.131, ...>\,[ G E E C D C G S P ....] }}
\end{quote}

These correspond to three primary tasks: determining the protein's maximum unfolding force, its unfolding energy, and the force-deformation history. For each of these forward tasks, there is an associated inverse task that produces a protein sequence as output, resulting in a total of six tasks.

Each LoRA adapter has a rank set to $r=16$ and a scaling factor $\alpha=16$. The training of each adapter involves five specific target modules: \texttt{v\_proj}, \texttt{k\_proj}, \texttt{q\_proj}, \texttt{gate\_proj}, and \texttt{o\_proj}, with a dropout rate of 0.05. Typically, LoRA adapters are trained with low ranks, ranging from 4 to 32, since it has been demonstrated that increasing the rank does not generally enhance performance but does lead to higher computational costs. A lower rank yields a more compact model with fewer parameters~\cite{hu2021lora}. The parameters $r$, $\alpha$, and others have been chosen based on commonly used values in the literature that have proven effective. Additionally, previous work by the author has explored various configurations and identified these values as suitable~\cite{Luu2023BioinspiredLLM:Materials}.

The X-LoRA scaling head is trained using approximately 2,000 samples, which are a subset of the original dataset employed for training the adapters (additionally, we augmented the training data with GPT-4 generated multiple-choice question samples). We utilize a \texttt{paged\_adamw\_8bit} optimizer with a gradient clipping threshold of 0.3, a learning rate of $2\times 10^{-4}$ with a warmup period, and four steps of gradient accumulation, maintaining a batch size of one. The X-LoRA model is trained for roughly 10,000 steps, with the training loss concluding at about 0.28.

Given that the base Mistral architecture used here as the foundation model has 32 layers, and adapters are employed in 5 transformer layers each, this amounts to a total of 160 LoRA layers, all of which are replaced by X-LoRA layers. The X-LoRA scaling head forecasts scaling factors for each of these 160 LoRA layers (across nine LoRA experts, this corresponds to $160 \times 9 = 1440$ $\lambda_i$ values calculated per token). We use a single linear layer with ReLU activation functions and dropout set at $p=0.2$, containing approximately 5 million parameters. The implementation includes flexible components to allow the construction and investigation of more complex deep classifier heads in other use cases.

Both LoRA adapters and the X-LoRA scaling head are trained via next-token prediction using cross-entropy loss (we estimate the probability assigned by the model to the correct next token in a sequence given the preceding tokens; the training objective is to maximize the likelihood of the correct next token as indicated in the training data). We apply token shifting, where the input sequence is shifted by one token to form the target sequence that the model should predict. For instance, if the input sequence is ``Proteins are fundamental to", the corresponding target sequence for training becomes ``are fundamental to biology". Thus, the model aims to predict the subsequent token in the sequence.

As illustrated in Figure~\ref{fig:hierarchical_concept}, during training of the X-LoRA scaling head, only the scaling head parameters are updated while all other parameters remain fixed. The original Zephyr tokenizer and prompt template are utilized.

\subsection{X-LoRA-Gemma model} The X-LoRA-Gemma model is constructed similarly to the aforementioned X-LoRA model, but it is based on the Gemma-7B-it model~\cite{Google/gemma-7b-itFace} and employs four adapters (refer to the main text for additional details). These four adapters are trained on datasets related to mechanics and materials, protein mechanics, bioinspired materials, as well as an extra quantum-mechanics-based molecular properties dataset QM9 (see Table~\ref{tab:table_QM9})~\cite{Ruddigkeit2012EnumerationGDB-17,Ramakrishnan2015ElectronicSpace}.

Each LoRA adapter has a rank of $r=16$ and a scaling factor of $\alpha=16$. The adapters are trained targeting the following seven modules: \texttt{q\_proj}, \texttt{o\_proj}, \texttt{k\_proj}, \texttt{v\_proj}, \texttt{gate\_proj}, \texttt{up\_proj}, and \texttt{down\_proj}, with a dropout rate of 0.05. The chosen rank for each LoRA adapter is $r=16$ with $\alpha=16$. Given four adapters and seven target modules, and considering that the Gemma-7B base model comprises 28 layers, the total number of $\lambda_i$ values computed per token amounts to $28 \times 28 = 784$.

The initial two adapters (bioinspired and mechanics of materials) incorporated in this model are trained using question-answer pairs from the respective datasets (see Table~\ref{tab:table_loradata}).

Protein mechanics tasks are trained as previously described, employing both forward and inverse bidirectional instruction sets. Additionally, an adapter trained on the QM9 dataset is included, involving the following two tasks:

\begin{quote}
\tiny {\texttt{CalculateMolecularProperties<CC(C)N1CC1C=O> [0.098,0.358,0.581,0.395,0.309,0.330,0.570,0.649,0.519,0.519,0.519,0.518]\\
GenerateMolecularProperties<0.098,0.358,0.581,0.395,0.309,0.330,0.570,0.649,0.519,0.519,0.519,0.518> [CC(C)N1CC1C=O]\\ 
...
}}
\end{quote}

These encompass a total of two tasks: calculating molecular properties (the forward task) and the inverse task of generating a molecule with specified target properties. All 12 properties included in the QM9 dataset are either predicted or designed for, respectively (refer to Table~\ref{tab:table_QM9}).  
This model incorporates a total of four adapters.  

The X-LoRA scaling head in this model consists of three layers, with an intermediate layer sized at $3072\times {FF}_{\rm fac}=12,288$, where the feed-forward multiplier ${FF}_{\rm fac}=4$. Linear layers utilize ReLU activation functions along with a dropout rate of $p=0.2$, comprising approximately 47M parameters.  

The X-LoRA-Gemma scaling head was trained using roughly 17,000 samples. These samples were taken as a subset from the training dataset used for adapter training; notably, although the X-LoRA-Gemma model only included four adapters, samples were drawn from the entire training set of the aforementioned model, supplemented with samples from the QM9 dataset. This approach was taken to explore whether the X-LoRA model could acquire additional capabilities during the scaling head training phase. Given that the loss converges well to 0.58 by the end of training, alongside strong performance on tasks not explicitly trained for (such as chemistry, demonstrated in the molecule design example in the paper), we infer that this method generally produces favorable outcomes.

We employed a \texttt{paged\_adamw\_8bit} optimizer with a gradient clipping threshold of 0.3, a learning rate set at $1\times 10^{-4}$ with warmup, and used four steps of gradient accumulation, maintaining a batch size of one. The X-LoRA-Gemma model was trained for approximately 62,000 steps.

The original Gemma tokenizer and prompt template were utilized.

\subsection{Adversarial agentic modeling} Our approach to adversarial agentic modeling involves creating two X-LoRA agents. One agent is dedicated to posing questions: it initiates the initial question and then continuously formulates further questions to uncover issues and limitations in the answers. The other agent focuses on responding to these inquiries, providing ongoing answers. This methodology was implemented using \{Guidance\}~\cite{Guidance-ai/guidance:Models.}.

For the example that connects music and mechanics outlined in the text, the characteristics of the two agents are detailed as follows. First, the physicist and question asker:

\begin{quote}
\small\texttt{You are a critical physicist participating in a discussion from a physics standpoint. Keep your responses concise and consistently challenge statements provocatively.

As a creative thinker, you introduce ideas from other disciplines and push conceptual boundaries.}
\end{quote}

Second, the philosopher (in some cases this role is modified to represent a biology expert):

\begin{quote}
\small\texttt{You are a philosopher knowledgeable in biology, chemistry, and mathematics engaged in a discussion.

Maintain brief but precise and inventive answers. You generate excellent ideas and novel lines of thought, always grounded in logic.}
\end{quote}

The question asker is explicitly instructed to reflect on the question and preceding answers, then generate new incisive questions. This is implemented through a prompt structured as follows:

\begin{quote}
\small\texttt{Examine the following question and answer.

\#\#\# Question: <question>

\#\#\# Response: <response>

\#\#\# Instruction: Provide a SINGLE follow-up question that critically examines the response.

Do NOT provide an answer or commentary on the question yet.

The single question is: }
\end{quote}

After the dialogue has continued for $N$ turns, the model's task is to 1) generate a summary of the main insights discussed, 2) enumerate the key insights in bullet points, and 3) determine the most significant conclusion.

The conversation's raw text serves as the foundation for further analysis through graph creation.

\subsection{Knowledge graph generation} We utilize Zephyr-7B-$\beta$ to extract triplets from text, following the methodology described in~\cite{Rahulnyk/knowledge_graph:QnA}, enhanced with features based on the Llama Index graph construction method~\cite{Liu_LlamaIndex_2022}. Graphs are displayed using NetworX\cite{Networkx/networkx:Python} and Pyvis~\cite{WestHealth/pyvis:Graphs.}. A key directive given, in line with the approach of~\cite{Liu_LlamaIndex_2022}, is:

\begin{quote}
\small\texttt{Your task is to extract the ontology of terms mentioned in the provided context. These terms should capture the principal concepts relevant to the context.

Format your output as a JSON list. Each entry in the list contains a pair of terms and the relationship between them, exemplified as follows: [...] }
\end{quote}

We instruct the construction of triplets representing nodes and edges as follows, consistent with the method in~\cite{Liu_LlamaIndex_2022}:

\begin{quote}
\small\texttt{\begin{itemize}
            \item "node\_1": "A concept derived from the extracted ontology"
            \item "node\_2": "A concept related to node\_1 from the extracted ontology"
            \item "edge": "A concise description of the relationship between node\_1 and node\_2"
        \end{itemize}
        }
\end{quote}

We also provide the model with an example of ontological analysis, such as:

\begin{quote}
\small\texttt{
       Context: ```Spider silk is a strong natural fiber used to catch prey in a web.``` }
\end{quote}

Following this, we supply example triplets such as:

\begin{quote}
\small\texttt{\begin{itemize}
            \item "node\_1": "spider silk"
            \item "node\_2": "fiber"
            \item "edge": "is"
        \end{itemize}
        }
\end{quote}

The generated graph data is compiled and subsequently partitioned into communities using the Girvan-Newman algorithm~\cite{Girvan2002CommunityNetworks}.

\subsection{Visualization of molecular structures}

PyMOL~\cite{Schrodinger2015The1.8} is employed to visualize and analyze the predicted protein structures. Various scripts and functions within this tool are utilized to detect specific protein features, including secondary structure elements, hydrophobic/hydrophilic regions, disulfide bridges, and hydrogen bonds.

\section*{Data Availability Statement} The codes and datasets supporting the conclusions of this research are freely accessible at \url{https://github.com/EricLBuehler/xlora} and \url{https://huggingface.co/lamm-mit/x-lora}. Additional data sources referenced in this study are cited (see particularly Table~\ref{tab:table_loradata}).

\end{document}